
\chapter{Diagrama de Clases Conceptuales}\label{ch:diagramaclases}

El diagrama de clases conceptuales representa la estructura estática del \emph{dominio del problema}. A diferencia de un diagrama de clases de software o de diseño, aquí no se incluyen tipos de datos, métodos, ni elementos de la capa de servicios, infraestructura o tecnología: solo conceptos del dominio, sus atributos esenciales y las asociaciones entre ellos. Estas clases conceptuales son las candidatas a convertirse, durante el diseño, en clases de software.

Para FisioManager se modelan las clases conceptuales agrupadas en paquetes que reflejan los \textit{bounded contexts}: Usuarios, Clínico, Agenda, Rutinas, Seguimiento, Comunicación y Plataforma.

\begin{figure}[H]
  \centering
  \includegraphics[width=\textwidth,keepaspectratio]{FisioManager_DiagramaClases.png}
  \caption{Diagrama de Clases Conceptuales del sistema FisioManager.}
  \label{fig:diagrama-clases}
\end{figure}

\section{Decisiones de modelado}

Las siguientes decisiones de modelado fueron tomadas durante el análisis orientado a objetos y se reflejan tanto en las CRC Cards como en el diagrama de clases conceptuales:

\begin{itemize}
  \item \textbf{Herencia en Usuario:} se modela una clase abstracta \codeinline{Usuario} de la cual heredan \codeinline{Paciente}, \codeinline{Terapeuta} y \codeinline{Administrador}. Esto centraliza los atributos comunes y evita duplicación a nivel conceptual.

  \item \textbf{Separación Rutina vs AsignacionRutina:} una \codeinline{Rutina} es una plantilla reutilizable; \codeinline{AsignacionRutina} representa la instancia concreta para un paciente con fechas y estado. Esto permite asignar la misma \codeinline{Rutina} a múltiples pacientes sin duplicarla.

  \item \textbf{PlantillaActividad vs Actividad:} el catálogo reutilizable se modela como \codeinline{PlantillaActividad}; cada \codeinline{Actividad} es la instancia concreta dentro de una \codeinline{Rutina}, con sus propios parámetros (duración, repeticiones, descanso y orden) que dependen del contexto de la rutina.

  \item \textbf{SesionEjecucion como agregado de seguimiento:} cada vez que el paciente realiza una rutina se crea una \codeinline{SesionEjecucion}. Esta es la entidad sobre la que pivotan el feedback y el cálculo de adherencia.

  \item \textbf{Feedback con resultados de IA en línea:} el resultado del procesamiento (transcripción de Whisper y resumen clínico de Ollama) se almacena como atributos del propio \codeinline{Feedback} (\codeinline{transcripcion}, \codeinline{resumenIA}); son opcionales porque solo existen cuando el paciente graba audio. No se modela una entidad separada de análisis.

  \item \textbf{Mapa corporal y escalas:} el \codeinline{PuntoDolor} registra la localización del dolor asociada a un \codeinline{Feedback}, y la \codeinline{EscalaResultado} captura mediciones clínicas objetivas (ROM, fuerza, escalas funcionales, VAS) que registra el terapeuta para seguir la evolución.

  \item \textbf{PlanTratamiento y FaseTratamiento:} el plan agrupa fases secuenciales y prescribe las rutinas que el paciente debe ejecutar. La separación entre \codeinline{PlanTratamiento} y \codeinline{AsignacionRutina} permite que el plan exprese intención clínica (objetivos, fases) mientras que la asignación expresa la prescripción operativa.

  \item \textbf{NotaClinica como concepto propio:} las notas que el terapeuta registra durante consultas no se modelan como simples atributos de texto en \codeinline{HistorialClinico} o \codeinline{PlanTratamiento}, sino como un concepto independiente con identidad propia, fecha y autor. Esto permite trazabilidad de la evolución clínica y reutilizar la nota en distintos contextos.

  \item \textbf{Paquete Agenda:} se introduce \codeinline{Cita} como concepto central de la planeación clínica. La cita ocurre entre un paciente y un terapeuta, opcionalmente ocupa una \codeinline{Sala} y puede ligarse a un \codeinline{PlanTratamiento}. El \codeinline{HorarioDisponible} expresa la oferta semanal del terapeuta y es la base contra la que se validan tanto las citas que crea el administrador como las que el paciente agenda por sí mismo. La \codeinline{ListaEspera} permite avisar al paciente cuando se libera un hueco.

  \item \textbf{Paquete Comunicación:} la mensajería en tiempo real se modela con \codeinline{Conversacion} (única entre un paciente y su terapeuta) que agrupa \codeinline{Mensaje}s. La \codeinline{Notificacion} representa el centro de avisos in-app y la \codeinline{PreferenciaNotificacion} guarda las preferencias de cada usuario (correo, recordatorios, horario silencioso).

  \item \textbf{Paquete Plataforma:} el \codeinline{RegistroAuditoria} captura las acciones sensibles del sistema (inicios de sesión, consulta de historial, gestión de documentos) con su actor, entidad afectada e IP, dando trazabilidad de cumplimiento.

  \item \textbf{DocumentoClinico:} los archivos del expediente (radiografías, laboratorios, consentimientos) se modelan como concepto propio del paquete Clínico, con su categoría y bandera de visibilidad para el paciente.
\end{itemize}

\section{Código PlantUML del diagrama}

El diagrama es totalmente compilable. El código fuente se entrega como archivo \codeinline{.puml} en el repositorio y puede regenerarse con el comando:

\begin{lstlisting}[language=bash]
plantuml -tpng 01_diagrama_clases.puml
\end{lstlisting}

El código fuente completo del diagrama se incluye en el Anexo~\ref{ch:anexos} al final de este documento.
